% Nouveau EGA V (cote 25): archive cote 25, site page 132 (PDF page 133).
% typescript. Draft compared: none.
% AI-transcribed by Claude Opus 5.5 (claude-opus-5-5, adaptive thinking, effort high), Message Batches API, batch ega25b (msgbatch_01A7n5YuXubUPNK1zsepi8RG), 2026-11-13; not yet checked by eye.
\scanpage{25/132}
\mspage{\hand{35}}

satisfaisant certaines conditions, cf \hand{4}.4.). D'autre part, \hand{4}.3. b)
équivaut \var{\xout{xxxxxtixltxxxxt}}{}{typewriter strike-out}{ega25_25_132-v0001} aussi à dire que \var{\inserted[typed]{lorsque g est surjectif,}}{lorsque g est surjectif,}{insertion}{ega25_25_132-v0002} si deux systêmes lin de diviseurs
sur X/S \var{\xout{xxxx}}{}{typewriter strike-out}{ega25_25_132-v0003} paramétrés par deux fibrés projectifs $\overset{\hand{\vee}}{P},\overset{\hand{\vee}}{P}'$ sont isomorphes;
alorsilexiste un \uline{unique} isomorphisme de $\overset{\hand{\vee}}{P}$ sur $\overset{\hand{\vee}}{P}$\unsure{(|'} compatibles avec D,D';
on peut donc dire qu'une classe (à isom près) de systêmes linéaires sur
X/S détermine son fibré projectif paramétrisateur \uline{à isomorphisme} \uuline{unique}
\uline{près} . Techniquement, ce r ésultat permettra (une fois qu'on disposera
de la te\unsure{ç}hnique de descente, exposée au Chap V) de faire de \uline{la d escente}
\uline{fid plate pour des systêmes linéaires de diviseurs}, - sous réserve
\var{\margin{? (}}{\placedmargin{? (}}{marginal note placed}{ega25_25_132-v0004} toutefois d'admettre \var{\inserted{également}}{également}{insertion}{ega25_25_132-v0005} comme fibrés paramétrisateurs des ``fibrés projectifs
tordus'', ce qui va être fait au prochain \var{\deleted{paragraphex}}{}{deletion}{ega25_25_132-v0006} numéro.

Descendan\var{\deleted{t}}{}{deletion}{ega25_25_132-v0007} à nouveau par terre, et bien bas, pour expliciter en
termes vulgaires la notion de système linéaire, en nous plaçant pour
simplifier sur un corps de base (bien que l'énoncé subsiste essentielle-%
ment tel quel sur une base affine):
\uline{Proposition} \var{\deleted{\hand{4}.4.}}{}{deletion}{ega25_25_132-v0008} Soit X un préschéma \var{\xout{\unsure{loc}}}{}{typewriter strike-out}{ega25_25_132-v0009} \var{\inserted[typed]{\var{\xout{xxxxxide}}{}{typewriter strike-out}{ega25_25_132-v0011}}}{\var{\xout{xxxxxide}}{}{typewriter strike-out}{ega25_25_132-v0011}}{insertion}{ega25_25_132-v0010} de type fini sur un corps k,
tel que $k\hand{\to} H^0(X,\underline{O}_X)$ soit un isomorphisme\var{\over{,}{.}}{.}{writing over}{ega25_25_132-v0012} A tout système linéaire \var{\inserted[typed]{D}}{D}{insertion}{ega25_25_132-v0013} de
diviseurs \var{\deleted{Dx}}{}{deletion}{ega25_25_132-v0014} \var{\inserted[typed]{sur X,}}{sur X,}{insertion}{ega25_25_132-v0015} paramétré par un fibré projectif $\overset{\hand{\vee}}{P}$ sur k, associons
l'ensemble\hand{(!)}Ens(D) de tous les diviseurs sur X de la forme D(t), où
$t\hand{\in}\overset{\hand{\vee}}{P}(k)$. \hand{a)} Si \var{\deleted{alors}}{}{deletion}{ega25_25_132-v0016} D' est un autre systême \var{\inserted[typed]{lin}}{lin}{insertion}{ega25_25_132-v0017} de diviseurs\var{\inserted[typed]{sur X}}{sur X}{insertion}{ega25_25_132-v0018}, paramétré par un
fibré projectif $\overset{\hand{\vee}}{D}'$ , alors D et D' sont isomorphes sss Ens(D)=Ens(D').
\hand{b)} \var{\inserted{Supposons k alg. clos et X \unsure{gén.\textsuperscript{t}|génér.\textsuperscript{t}} intègre.}}{Supposons k alg. clos et X \unsure{gén.\textsuperscript{t}|génér.\textsuperscript{t}} intègre.}{insertion}{ega25_25_132-v0019}
\hand{(}Pour qu'un ensemble $\hand{\Delta}$ de diviseurs de Cartier positifs sur X soit de
le forme Ens(D), il f et s qu'il \var{\inserted{existe}}{existe}{insertion}{ega25_25_132-v0020} \var{\deleted{soit vide, ou qu'il satisfasse la condi-%
tion suivante: \var{\xout{pourxxxxsoitxxD$_0$xxxxxxxxxalors}}{}{typewriter strike-out}{ega25_25_132-v0022}. il existe \var{\xout{desxfonctions}}{}{typewriter strike-out}{ega25_25_132-v0023}
\var{\xout{rationnellesxrelxàxkxxfxxx$\phi_1$xx}}{}{typewriter strike-out}{ega25_25_132-v0024} une suite finie $(\phi_i)_{1\ \ i\ \ n}$ de fonctions
rationelles rel k sur X, tellex que pour tout système d'eléments
$(c_i)_{1\ \ i\ \ n}$ de k}}{}{deletion}{ega25_25_132-v0021}

\mspage{\hand{132}}
